\Article{10}
Später versuchte auch der berühmte \textsc{de Foncenex}, nachdem er in \textsc{Eulers} erstem Beweis einen Mangel bemerkt hatte (oben Artikel 8, Einwand 4), den er nicht beheben konnte, noch einen anderen Weg und legte ihn in der genannten Abhandlung auf S. 120 vor\HistoricalNote{7}. Er besteht in Folgendem.

Es sei die Gleichung $Z=0$ vorgelegt, wobei $Z$ eine Funktion vom Grad $m$ der Unbekannten $z$ bezeichnet. Ist $m$ eine ungerade Zahl, so steht bereits fest, dass diese Gleichung eine reelle Wurzel hat. Ist aber $m$ gerade, so versucht der berühmte Foncenex auf folgende Weise zu beweisen, dass die Gleichung wenigstens eine Wurzel der Form $p+q\BareRoot{-1}$ hat. Es sei $m=2^n i$, wobei $i$ eine ungerade Zahl bezeichnet, und es werde vorausgesetzt, dass $zz+uz+M$ ein Teiler der Funktion $Z$ ist. Dann sind die einzelnen Werte von $u$ die Summen von je zwei Wurzeln der Gleichung $Z=0$ mit umgekehrtem Vorzeichen; daher hat $u$ insgesamt $\frac{m\mathbin{.}m-1}{1\mathbin{.}2}=m'$ Werte. Wird vorausgesetzt, dass $u$ durch eine Gleichung $U=0$ bestimmt wird (wobei $U$ ein Polynom in $u$ und den bekannten Koeffizienten in $Z$ bezeichnet), so hat diese den Grad $m'$. Man sieht aber leicht, dass $m'$ eine Zahl der Form $2^{n-1}i'$ sein wird, wobei $i'$ eine ungerade Zahl bezeichnet. Ist nun $m'$ nicht ungerade, so werde wiederum $uu+u'u+M'$ als Teiler von $U$ vorausgesetzt. Durch ähnliche Überlegungen wird klar, dass $u'$ durch eine Gleichung $U'=0$ bestimmt wird, wobei $U'$ eine Funktion vom Grad $\frac{m'\mathbin{.}m'-1}{1\mathbin{.}2}$ in $u'$ ist. Setzt man $\frac{m'\mathbin{.}m'-1}{1\mathbin{.}2}=m''$, so ist $m''$ eine Zahl der Form $2^{n-2}i''$, wobei $i''$ eine ungerade Zahl bezeichnet. Ist nun $m''$ nicht ungerade, so werde $u'u'+u''u'+M''$ als Teiler der Funktion $U'$ angenommen. Dann wird $u''$ durch eine Gleichung $U''=0$ bestimmt, die, wenn man ihren Grad gleich $m'''$ setzt, für $m'''$ eine Zahl der Form $2^{n-3}i'''$ ergibt. Offenbar ist in der Folge der Gleichungen $U=0$, $U'=0$, $U''=0$ usw. die Gleichung an der Stelle $n$ von ungeradem Grad und hat daher eine reelle Wurzel. Wir setzen der Kürze halber $n=3$, sodass die Gleichung $U''=0$ eine reelle Wurzel $u''$ hat; denn man sieht ohne Mühe, dass dieselbe Schlussweise für jeden anderen Wert von $n$ gilt. Dann behauptet der berühmte de Foncenex, dass sich der Koeffizient $M''$ durch $u''$ und die Koeffizienten in $U'$ (die, wie man leicht erkennt, Polynome in den Koeffizienten in $Z$ sind), also durch $u''$ und die Koeffizienten in
\SourceLeaf{024}
$Z$ rational bestimmen lässt und daher reell ist. Hieraus folgt, dass die Wurzeln der Gleichung $u'u'+u''u'+M''=0$ in der Form $p+q\BareRoot{-1}$ enthalten sein werden; sie werden aber offenbar zugleich die Gleichung $U'=0$ erfüllen. Es gibt daher einen Wert von $u'$, der in der Form $p+q\BareRoot{-1}$ enthalten ist. Nun lässt sich der Koeffizient $M'$ auf dieselbe Weise wie zuvor rational durch $u'$ und die Koeffizienten in $Z$ bestimmen und wird daher ebenfalls in der Form $p+q\BareRoot{-1}$ enthalten sein. Somit werden die Wurzeln der Gleichung $uu+u'u+M'$ in derselben Form enthalten sein und zugleich die Gleichung $U=0$ erfüllen; das heißt, diese Gleichung wird eine Wurzel in der Form $p+q\BareRoot{-1}$ haben. Schließlich folgt hieraus auf ähnliche Weise, dass auch $M$ in derselben Form enthalten ist, ebenso eine Wurzel der Gleichung $zz+uz+M=0$, die offenbar auch die vorgelegte Gleichung $Z=0$ erfüllen wird. Daher wird jede Gleichung wenigstens eine Wurzel der Form $p+q\BareRoot{-1}$ haben.

\Article{11}
Die Einwände 1, 2 und 3, die ich gegen \textsc{Eulers} ersten Beweis erhoben habe (Artikel 8), haben gegen diese Methode dieselbe Beweiskraft, jedoch mit dem Unterschied, dass der zweite Einwand, dem \textsc{Eulers} Beweis nur in einigen besonderen Fällen ausgesetzt war, den vorliegenden in allen Fällen treffen muss. Man kann nämlich a priori beweisen: Selbst wenn eine Formel gegeben ist, die den Koeffizienten $M'$ rational durch $u'$ und die Koeffizienten in $Z$ ausdrückt, muss diese für mehrere Werte von $u'$ notwendig unbestimmt werden. Ebenso wird die Formel, die den Koeffizienten $M''$ durch $u''$ darstellt, für gewisse Werte von $u''$ unbestimmt usw. Dies erkennt man am deutlichsten, wenn man eine Gleichung vierten Grades als Beispiel nimmt. Wir setzen also $m=4$, und die Wurzeln der Gleichung $Z=0$ seien $\alpha$, $\beta$, $\gamma$, $\delta$. Dann ist klar, dass die Gleichung $U=0$ sechsten Grades sein wird und ihre Wurzeln $-(\alpha+\beta)$, $-(\alpha+\gamma)$, $-(\alpha+\delta)$, $-(\beta+\gamma)$, $-(\beta+\delta)$, $-(\gamma+\delta)$ sind. Die Gleichung $U'=0$ ist dagegen fünfzehnten Grades, und die Werte von $u'$ sind
\[
\begin{gathered}
2\alpha+\beta+\gamma,\quad 2\alpha+\beta+\delta,\quad 2\alpha+\gamma+\delta,\\
2\beta+\alpha+\gamma,\quad 2\beta+\alpha+\delta,\quad 2\beta+\gamma+\delta,\\
2\gamma+\alpha+\beta,\quad 2\gamma+\alpha+\delta,\quad 2\gamma+\beta+\delta,\\
2\delta+\alpha+\beta,\quad 2\delta+\alpha+\gamma,\quad 2\delta+\beta+\gamma,\\
\alpha+\beta+\gamma+\delta,\quad\alpha+\beta+\gamma+\delta,\quad\alpha+\beta+\gamma+\delta
\end{gathered}
\]
Bei dieser Gleichung muss man nun stehen bleiben, da ihr Grad ungerade ist. Sie wird tatsächlich die reelle Wurzel $\alpha+\beta+\gamma+\delta$ haben (die dem ersten Koeffizienten in $Z$ mit umgekehrtem Vorzeichen gleich und daher nicht nur reell, sondern auch rational sein wird, wenn die Koeffi\SourceLeaf{025}zienten in $Z$ rational sind). Man kann aber ohne Mühe erkennen: Ist eine Formel gegeben, die den Wert von $M'$ rational durch den zugehörigen Wert von $u'$ ausdrückt, so wird sie notwendig für $u'=\alpha+\beta+\gamma+\delta$ unbestimmt. Denn dieser Wert ist \textit{dreimal} eine Wurzel der Gleichung $U'=0$, und ihm entsprechen drei Werte von $M'$, nämlich $(\alpha+\beta)(\gamma+\delta)$, $(\alpha+\gamma)(\beta+\delta)$ und $(\alpha+\delta)(\beta+\gamma)$, die alle irrational sein können. Offenbar könnte aber eine rationale Formel in diesem Fall weder einen irrationalen Wert von $M'$ noch drei verschiedene Werte liefern. Aus diesem Beispiel lässt sich hinreichend entnehmen, dass die Methode des berühmten \textsc{de Foncenex} keineswegs befriedigt. Soll sie aber in jeder Hinsicht vollständig werden, so muss die Eliminationstheorie weit tiefer untersucht werden.

\Article{12}
Schließlich hat der ausgezeichnete \textsc{La Grange} unseren Satz in der Abhandlung \textit{Sur la forme des racines imaginaires des équations, Nouv. Mém. de l'Acad. de Berlin} 1772, S. 222 ff., behandelt. Dieser große Mathematiker hat sich besonders bemüht, die Mängel in \textsc{Eulers} erstem Beweis zu ergänzen, und tatsächlich vor allem das, was oben (Artikel 8) den zweiten und vierten Einwand ausmacht, so tief untersucht, dass nichts weiter zu wünschen bleibt, außer dass vielleicht in der früheren Untersuchung über die Eliminationstheorie, auf der diese ganze Untersuchung beruht, einige Zweifel übrig zu sein scheinen. --- Den dritten Einwand hat er jedoch überhaupt nicht berührt; vielmehr baut die ganze Untersuchung auf der Voraussetzung auf, jede Gleichung vom Grad $m$ habe tatsächlich $m$ Wurzeln.

Nach sorgfältiger Erwägung dessen, was bisher dargelegt wurde, hoffe ich daher, dass den Sachkundigen ein neuer, aus ganz anderen Grundsätzen abgeleiteter Beweis dieses höchst wichtigen Satzes nicht unwillkommen sein wird. Diesen darzustellen, gehe ich nun sogleich an.

\Article{13}
\textsc{Lemma.} \textit{Bezeichnet $m$ eine beliebige positive ganze Zahl, so ist die Funktion $\sin\varphi\mathbin{.}x^m-\sin m\varphi\mathbin{.}r^{m-1}x+\sin(m-1)\varphi\mathbin{.}r^m$ durch $xx-2\cos\varphi\mathbin{.}rx+rr$ teilbar.}

\textit{Beweis.} Für $m=1$ wird jene Funktion $=0$ und ist daher durch jeden Faktor teilbar; für $m=2$ wird der Quotient gleich $\sin\varphi$, und für jeden größeren Wert ist der Quotient $\sin\varphi\mathbin{.}x^{m-2}+\sin2\varphi\mathbin{.}rx^{m-3}+\sin3\varphi\mathbin{.}rrx^{m-4}+\text{usw.}+\sin(m-1)\varphi\mathbin{.}r^{m-2}$. Denn man bestätigt leicht, dass nach Multiplikation dieser Funktion mit $xx-2\cos\varphi\mathbin{.}rx+rr$ das Produkt der vorgelegten Funktion gleich wird.
